%----------------------------------------------------------------------------------------
%	PACKAGES AND OTHER DOCUMENT CONFIGURATIONS
%----------------------------------------------------------------------------------------
% Undergrad GPA 3.5/4.0 worth to put in?

% Section sequence?

% point 10 too small? Should I make it to 11 and make my resume to 3 pages?

% Do I need bullet points for teaching experience and make my resume to 3 pages?

% I actually have one more internship experience at Metagames when I was in undergrad, but I did not put it because the role was the same what I did in Bizbook and I did not want to make my CV to be more than 2 pages. Should I put it in?

% Should I include presentations?

% How can I improve my current CV?

\documentclass{resume} % Use the custom resume.cls style

\usepackage[left=0.75in,top=0.6in,right=0.75in,bottom=0.6in]{geometry} % Document margins
\usepackage{fontawesome}
\usepackage{hyperref}
\usepackage{enumitem}
\newcommand{\tab}[1]{\hspace{.2667\textwidth}\rlap{#1}}
\newcommand{\itab}[1]{\hspace{0em}\rlap{#1}}

\name{MYEONGSOO KIM}
\address{\faMapMarker~Santa Clara, CA 95054 \quad \faPhone~(+1) 678-576-5708 \quad \faEnvelope~codingsoo93@gmail.com \quad \faHome~\href{https://codingsoo.github.io/}{https://codingsoo.github.io/}}

\hypersetup{
    colorlinks=true,       % Enable colored links
    linkcolor=black,      % Color of internal links
    citecolor=black,      % Color of citations
    filecolor=black,      % Color of file links
    urlcolor=black        % Color of external links
}
\setlist[itemize]{topsep=-3pt, itemsep=0pt, partopsep=0pt, parsep=0pt}


\begin{document}

%----------------------------------------------------------------------------------------
%	SUMMARY SECTION
%----------------------------------------------------------------------------------------

\begin{rSection}{Summary}
Applied Scientist working on coding agents, evaluation design, and LLM-guided software engineering. Designed structured action spaces that reduce token consumption by up to 38\% while improving pass@1 accuracy (ACL 2026), built multi-turn evaluation benchmarks adopted org-wide to gate agent updates (NeurIPS 2025), and fine-tuned small language models achieving state-of-the-art test coverage with 80\% cost reduction over baselines (FSE 2025). Earlier work applied reinforcement learning to train exploration policies for automated API testing (ASE 2023, ICSE 2025). Currently building a self-improvement loop at AWS that mines real user traffic for agent failures, reproduces them as benchmarks, and validates fixes through reproducibility and fairness gates. 16 publications at top venues (ICSE, FSE, ASE, ISSTA, NeurIPS, ACL). Interested in post-training and RL-based approaches to improving code agent capabilities.
\end{rSection}

%----------------------------------------------------------------------------------------
%	WORK EXPERIENCE SECTION
%----------------------------------------------------------------------------------------


\begin{rSection}{Work Experience}
{\bf Amazon Web Services}, California, United States \hfill {\em Dec 2024 – Present} \\
{\bf Applied Scientist II (L5)}, AWS AI Labs — Kiro Org
\begin{itemize}
    \item Designed CodeStruct, an AST-based structured action space for code agents that reduces token consumption by up to 38\% while improving pass@1 accuracy; filed patent and published at ACL 2026.
    \item Building Kiro Trace, a self-improvement loop that mines live user traffic for agent failures, reproduces them as runnable benchmarks, and validates fixes through reproducibility and fairness gates before shipping.
    \item Built CodeAssistBench (NeurIPS 2025), a multi-turn evaluation benchmark for chat-based code assistance, adopted org-wide to gate all prompt, tool, and routing updates.
    \item Led agent architecture improvements for Kiro --- planning workflows, tool orchestration, and prompt engineering --- reducing median response latency by 14\% across 1M+ production messages; two patents filed.
    \item Developed TrajEval, a diagnostic evaluation framework for ``Coherence Collapse'' --- diagnosing why code agents fail after reaching the right code by measuring per-capability (search, read, write) success via oracle trajectory grounding (under review, EMNLP 2026).
\end{itemize}
{\bf IBM Research}, New York, United States \hfill {\em May 2024 – Aug 2024} \\
{\bf PhD AI Research Internship} with the AI for Code Team
\begin{itemize}
    \item Automated Unit Test Generation Project: Co-developed static analysis-guided LLM test 
  generation system combining symbolic execution with LLM-based assertion generation. 
  Integrated retrieval-augmented generation for framework-specific patterns, leading to 
  ICSE 2025 submission.
  \item Python Unit Test Generator: Designed and trained the Python component using supervised 
  fine-tuning on 50k+ test examples, achieving 9.8\% improvement in line coverage, 26.5\% 
  in branch coverage, and 22.5\% in method coverage vs. Microsoft's CodaMOSA. Conducted 
  160-participant survey at IBM demonstrating superior effectiveness over baselines.
  \item Java Unit Test Generator: Evaluated Java component using mutation coverage and semantic 
  correctness metrics, benchmarking against EvoSuite across 200+ open-source projects to 
  assess test quality and fault-detection capability.
  \item Multi-Agent Integration Test Generator: Architected multi-agent system using hierarchical 
  planning where specialized agents collaborate on API endpoint discovery (static analysis 
  agent), dependency resolution (constraint solver agent), and test case synthesis (LLM 
  generation agent). Designed inter-agent communication protocols for REST API testing 
  orchestration.
\end{itemize}

{\bf Samsung}, Seoul, South Korea \hfill {\em December 2023 – January 2024} \\
{\bf PhD AI Internship} with the AI Team in the Fire \& Marine Insurance Department
\begin{itemize}
    \item Developed BoToPa, an LLM-based chatbot using multi-layer RAG (GPT-4 + custom BM25 tokenizer + KNN text search) to match insurance inquiries to documents; achieved 90\% cost reduction and 10\% accuracy improvement over the previous system.
\end{itemize}

{\bf Google}, California, United States \hfill {\em May 2022 – August 2022} \\
{\bf PhD SWE Internship} with the Google Chat team at Google Cloud
\begin{itemize}
    \item Analyzed potential security issues in chat applications.
    \item Developed a grammar-based testing tool for internal and external chat applications which results in finding a severe security issue in Google Chat application.
\end{itemize}

{\bf IBM Research}, New York, United States \hfill {\em June 2020 – November 2020} \\
{\bf PhD Research Internship} with the AI for Code Team at the Thomas J. Watson Research Center
\begin{itemize}
    \item Investigated challenges in decomposing legacy web applications using Kubernetes.
    \item Developed a Command Line Interface tool for validating transformation of the monolithic application to microservice application.
\end{itemize}

\end{rSection}

%----------------------------------------------------------------------------------------
%	EDUCATION SECTION
%----------------------------------------------------------------------------------------

\begin{rSection}{Education}
{\bf Georgia Institute of Technology}, Georgia, United States \hfill {\em Aug 2019 -- Dec 2024} \\
Ph.D. in Computer Science, Advisor: Prof. Alessandro Orso
\begin{itemize}
    \item Developed multi-agent reinforcement learning for REST API test generation, training collaborative agents with custom reward shaping to maximize code/fault coverage (ASE 2023, ISSTA 2022).
    \item Designed semantic graph-based state representation enabling agents to reason about resource dependencies across 15+ real-world services, achieving 34\% higher branch coverage than baselines (ICSE 2025).
    \item Built LLM-guided test generation combining fine-tuned small language models with reinforcement learning, reducing generation cost by 80\% while maintaining effectiveness (FSE 2025).
    \item Enhanced API documentation understanding through NLP techniques (BERT fine-tuning, semantic similarity, entity extraction), reducing invalid requests by 45\% (ISSTA 2023).
\end{itemize}

{\bf Kookmin University}, Seoul, South Korea \hfill {\em Mar 2015 -- Aug 2019} \\
B.S. in Computer Science
\end{rSection}


%----------------------------------------------------------------------------------------
%	PUBLICATIONS SECTION
%----------------------------------------------------------------------------------------

\begin{rSection}{Publications}
\vspace{0.5em}
\textit{Note: ICSE, FSE, ASE, and ISSTA are widely recognized as the top four conferences in the field of Software Engineering.}

\textbf{Myeongsoo Kim}, Dingmin Wang, Siwei Cui, Farima Farmahinifarahani, Terry Yue Zhuo, Shweta Garg, Baishakhi Ray, Rajdeep Mukherjee, Varun Kumar. \textit{Coherence Collapse: Diagnosing Why Code Agents Fail After Reaching the Right Code}. Under review at EMNLP 2026.

\textbf{Myeongsoo Kim}, Joe Hsu, Dingmin Wang, Shweta Garg, Varun Kumar, Murali Krishna Ramanathan. \textit{CodeStruct: Code Agents over Structured Action Spaces}. In Proceedings of the 64th Annual Meeting of the Association for Computational Linguistics (ACL). 2026.

\textbf{Myeongsoo Kim}, Shweta Garg, Baishakhi Ray, Varun Kumar, Anoop Deoras. \textit{CodeAssistBench (CAB): Dataset \& Benchmarking for Multi-turn Chat-Based Code Assistance}. In Proceedings of NeurIPS 2025, Datasets \& Benchmarks Track. 2025.

Tyler Stennett, \textbf{Myeongsoo Kim}, Saurabh Sinha, Alessandro Orso. \textit{AutoRestTest: A Tool for Automated REST API Testing Using LLMs and MARL}. In Proceedings of the 47th IEEE/ACM International Conference on Software Engineering (ICSE), Demonstrations Track. 2025.

\textbf{Myeongsoo Kim}, Saurabh Sinha, Alessandro Orso. \textit{LlamaRestTest: Effective REST API Testing with Small Language Models}. Proceedings of the ACM on Software Engineering 2.FSE (2025): 465-488.

Rangeet Pan, \textbf{Myeongsoo Kim}, Rahul Krishna, Raju Pavuluri, Saurabh Sinha. \textit{Aster: Natural and Multi-Language Unit Test Generation with LLMs}. In Proceedings of the 47th IEEE/ACM International Conference on Software Engineering (ICSE), Software Engineering in Practice (SEIP) Track. 2025. \textbf{Distinguished Paper Award}.

\textbf{Myeongsoo Kim}, Saurabh Sinha, Alessandro Orso. \textit{A Multi-Agent Approach for REST API Testing with Semantic Graphs and LLM-Driven Inputs}. In Proceedings of the 47th IEEE/ACM International Conference on Software Engineering (ICSE). 2025.

\textbf{Myeongsoo Kim}, Santosh Pande, Alessandro Orso. \textit{Improving Program Debloating with 1-DU Chain Minimality}. 2024 IEEE/ACM 46th International Conference on Software Engineering: Companion Proceedings (ICSE-Companion). p. 384–385.

\textbf{Myeongsoo Kim}, Tyler Stennett, Dhruv Shah, Saurabh Sinha, Alessandro Orso. \textit{Leveraging Large Language Models to Improve REST API Testing}. In Proceedings of the ICSE-NIER 2024, IEEE/ACM International Conference on Software Engineering - New Ideas and Emerging Results Track (ICSE-NIER). 2024. p. 37-41.

\textbf{Myeongsoo Kim}, Saurabh Sinha, Alessandro Orso. \textit{Adaptive REST API Testing with Reinforcement Learning}. In Proceedings of the 38th IEEE/ACM International Conference on Automated Software Engineering (ASE). 2023. p. 446-458.

\textbf{Myeongsoo Kim}, Davide Corradini, Michele Pasqua, Mariano Ceccato, Alessandro Orso, Saurabh Sinha, Rachel Tzoref-Brill, Mariano Ceccato. \textit{Enhancing REST API Testing with NLP Techniques}. In Proceedings of the 32nd ACM SIGSOFT International Symposium on Software Testing and Analysis (ISSTA). 2023. p. 1232-1243.

\textbf{Myeongsoo Kim}, Qi Xin, Saurabh Sinha, Alessandro Orso. \textit{Automated Test Generation for REST APIs: No Time to Rest Yet}. In Proceedings of the 31st ACM SIGSOFT International Symposium on Software Testing and Analysis (ISSTA). 2022. p. 289-301.

Qi Xin, \textbf{Myeongsoo Kim}, Qirun Zhang, Alessandro Orso. \textit{Subdomain-Based Generality-Aware Debloating}. In Proceedings of the 35th IEEE/ACM International Conference on Automated Software Engineering (ASE). 2020. p. 224-236.

Qi Xin, \textbf{Myeongsoo Kim}, Qirun Zhang, Alessandro Orso. \textit{Program Debloating via Stochastic Optimization}. In Proceedings of the ACM/IEEE 42nd International Conference on Software Engineering: New Ideas and Emerging Results (ICSE-NIER). 2020. p. 65-68.

\textbf{Myeongsoo Kim}, Changheon Song, Hyeji Kim, Deahyun Park, Yeeji Kwon, Eun Namkung, Ian G. Harris, Marcel Carlsson. \textit{Scam Detection Assistant: Automated Protection from Scammers}. In 2019 First International Conference on Societal Automation (SA). IEEE, 2019. p. 1-8.

Elahe Paikari, JaeEun Choi, SeonKyu Kim, Sooyoung Baek, \textbf{Myeongsoo Kim}, SeungEon Lee, ChaeYeon Han, YoungJae Kim, KaHye Ahn, Chan Cheong, André van der hoek. \textit{A chatbot for conflict detection and resolution}. In 2019 IEEE/ACM 1st International Workshop on Bots in Software Engineering (BotSE). IEEE, 2019. pp. 29-33.

\textbf{Myeongsoo Kim}, Eun-Jin Im. \textit{Ethereum Smart Contract Vulnerability Detector}. Journal of the Korea Information Science Association, 2018, pp. 1940-1942.

\end{rSection}


%----------------------------------------------------------------------------------------
%	TEACHING EXPERIENCE SECTION
%----------------------------------------------------------------------------------------

% \begin{rSection}{Teaching Experience}
% {\bf Georgia Institute of Technology}, GA, United States \hfill {\em August 2021 – December 2021} \\
% Teaching Assistant, Advanced Software Engineering (CS 6301)
% \end{rSection}


%----------------------------------------------------------------------------------------
%	PRESENTATIONS SECTION
%----------------------------------------------------------------------------------------




%----------------------------------------------------------------------------------------
%	AWARDS SECTION
%----------------------------------------------------------------------------------------

\begin{rSection}{Awards}

\textbf{Distinguished Paper Award}, ICSE 2025 (SEIP Track) \hfill {\em April 2025} \\
\textbf{Third Prize}, Hand in Hand Challenge \hfill {\em January 2021} \\
\textbf{CS 7001 Award}, GT Computing Student Awards \hfill {\em April 2020} \\
\textbf{Undergraduate Research Fellowship}, Kookmin University \hfill {\em 2017-2019}\\
\textbf{First Prize}, LINE Blockchain Competition \hfill {\em December 2018} \\
\textbf{Second Prize}, KMU Capstone Design Competition \hfill {\em November 2018} \\
\textbf{Best Undergraduate Paper Award}, Korea Computer Congress \hfill {\em July 2018} \\
\textbf{Software Creativity Award (First Prize)}, President of Korea Science Foundation \hfill {\em March 2017}
\end{rSection}

%----------------------------------------------------------------------------------------
%	TECHNICAL SKILLS SECTION
%----------------------------------------------------------------------------------------

% \begin{rSection}{Technical Skills}
% \textbf{Language}: Python, Java, C \\
% \textbf{OS}: Windows, Mac, Ubuntu, Debian, CentOS \\
% \textbf{Domain Experience}: Software Testing, REST API, Natural Language Processing, Neural Networks, Reinforcement Learning, Blockchain 
% \end{rSection}

%----------------------------------------------------------------------------------------
%	SERVICE AND LEADERSHIP SECTION
%----------------------------------------------------------------------------------------

% \begin{rSection}{Service and Leadership}
% Reviewer, ACM Transactions on Software Engineering and Methodology \hfill {\em 2025} \\
% Reviewer, IEEE ACCESS \hfill {\em 2025} \\
% Reviewer, Jordanian Journal of Computers and Information Technology \hfill {\em 2024} \\
% President, Georgia Tech Korean Student Association \hfill {\em 2020 - 2021} \\
% Web Developer, Georgia Tech Korean Student Association \hfill {\em 2019 - 2020}\\
% Volunteer for Making Braille Books for the Blind, Song-Pa Community Center \hfill {\em January 2017 - February 2017}\\
% Sergeant, Republic of Korea Army \hfill {\em August 2014 – May 2016} 
% \end{rSection}

\end{document}
